% LaTeX fragment: Chapter VI foundation block, 17 October 2026.
% Requires amsmath and amssymb; no preamble or theorem package is assumed.
% Agent-originated analytical successor; no native calculation is asserted.

\section{A quadratic overlap enclosure for the declared continuum preparation}
\label{sec:chvi-foundation-overlap}

\paragraph{Provenance and scope.}
The quadratic refinement and robustness corollaries in this section are
agent-originated derivations by OpenAI ChatGPT and Codex in the present work
block. They use the retained Chapter~VI Hamiltonian, exact hard-eigenbranch
preparation and continuum form construction. They do not replace that
preparation by bare source levels, change its hard--soft split, or assert a
new native certificate or a laboratory preparation. Route interpretation
continues to use the chapter's declared proper-time controlled evolution.

\subsection{Hamiltonian, preparation and operator-domain hypotheses}

Retain the complete sixteen-state source and both continuum polarizations:
\begin{equation}
\begin{split}
 H_0={}&\epsilon(S+Q_L+Q_R+Q_LQ_R)+H_f
       -\epsilon(D\phi_++C\phi_\times)\\
      &+\epsilon^2\mu(D^2+C^2),
 \\ &D=Q_L-Q_R,\quad C=iSD,\quad \mu=(4R)^{-1}.
\end{split}
\label{eq:chvi-foundation-full-generator}
\end{equation}
The result below concerns one invariant, separately prepared lane
$a<b$, with signed $d=q_a-q_b<0$, $g=\epsilon d$, and
$e=q_a+q_b+q_aq_b$. The restriction of this same generator is
\begin{equation}
 H_0=\epsilon eI+\epsilon\sigma_x+H_f
   -g(\sigma_z\phi_++\sigma_y\phi_\times)+2g^2\mu I.
\label{eq:chvi-foundation-lane-generator}
\end{equation}
The lane restriction specifies the preparation sector; it does not discard
the remaining source sectors from the model. No claim about an arbitrary
superposition of different lanes is made here.

Use the chapter's declared split $\lambda=64\Delta$ and factor the field
into hard and soft Fock factors, including both polarizations in each.
Let $H_{\rm ref}$ retain the hard interaction, free soft field, even source
energy and the \emph{same full counterterm} $2g^2\mu I$ as $H_0$.
Thus their closed quadratic forms have difference
\begin{equation}
\begin{gathered}
V_s=-g(\sigma_z\phi_{+,s}+\sigma_y\phi_{\times,s}),
 \\
 v_s=\int_0^\lambda |u_R(\omega)|^2\,\mathrm d\omega
 \leq v_s^{\rm up}
 =\frac{R^4\lambda^6}{675\kappa\pi}.
\end{gathered}
\label{eq:chvi-foundation-soft-difference}
\end{equation}
In particular, no counterterm is subtracted a second time. The signed $g$
remains in the interaction; its square below comes from a norm.

Both Hamiltonians are the self-adjoint, semibounded operators associated
with the retained closed forms on the common form domain
$\mathcal Q=D(H_f^{1/2})$, with the finite source factor understood.
The full ultraviolet smearing need not be square integrable. The soft
smearing in Eq.~\eqref{eq:chvi-foundation-soft-difference} is square
integrable, and the chapter's weighted field estimates define the common
forms. We do not treat $V_s$ as a bounded operator on the full Fock space
or subtract unbounded operators on an unspecified domain.

Let the normalized orthogonal exact hard eigenbranches, including the
matched scalar terms, give
\begin{equation}
\begin{gathered}
\psi_k=\Phi_k\otimes|0_s\rangle\in D(H_{\rm ref}),\\
 H_{\rm ref}\psi_k=E_k\psi_k,\\
 \Psi_0=\frac{\psi_0+\psi_1}{\sqrt2},\\
 \Omega=E_1-E_0>0.
\end{gathered}
\label{eq:chvi-foundation-preparation}
\end{equation}
Here $|0_s\rangle$ is the joint soft vacuum of both polarizations.
Set
\begin{equation}
\begin{gathered}
U(t)=e^{-itH_0},\\
 R(t)=e^{-itH_{\rm ref}},\\
 b^2=2g^2v_s,\\
 \bar E=\frac{E_0+E_1}{2}.
\end{gathered}
\label{eq:chvi-foundation-notation}
\end{equation}
The symbol $R(t)$ denotes the reference propagator, whereas the scalar
$R$ in the form factor is the retained source radius.

For every real $u$, the form action of $V_s$ on $R(u)\Psi_0$ is represented
by the Hilbert-space vector
$f(u)=V_sR(u)\Psi_0$. It lies in the one-soft-photon sector and is
continuous in $u$. The two polarization insertions are orthogonal, while
$\sigma_z$ and $\sigma_y$ preserve the source norm. Consequently
\begin{equation}
\begin{gathered}
\|f(u)\|^2=b^2,
 \\
 \langle\Psi_0,R(v)f(u)\rangle=0
 \quad(u,v\in\mathbb R).
\end{gathered}
\label{eq:chvi-foundation-forcing-orthogonality}
\end{equation}
The last identity uses the zero-soft-photon reference orbit, not a
truncation of the actual evolution.

For completeness, the representation theorem for the common closed form
gives, against every test vector in $\mathcal Q$,
\begin{equation}
\begin{gathered}
R(u)\Psi_0\in D(H_0),\\
 H_0R(u)\Psi_0=H_{\rm ref}R(u)\Psi_0+f(u).
\end{gathered}
\label{eq:chvi-foundation-domain}
\end{equation}
The right side is Hilbert-space continuous, since the reference orbit is
a finite sum of eigenvectors and $f$ is continuous. These are sufficient
strong-Duhamel hypotheses. At $u=0$ they also give
\begin{equation}
\begin{gathered}
\langle H_0\rangle_{\Psi_0}=\bar E,\\
 \sigma_H^2:=\|(H_0-\bar E)\Psi_0\|^2
       =\frac{\Omega^2}{4}+b^2.
\end{gathered}
\label{eq:chvi-foundation-variance}
\end{equation}
This uses $\Psi_0\in D(H_0)$; neither $D(H_0^2)$ nor a fourth energy
moment is assumed. No numerical value of $\bar E$ is supplied here.

\subsection{The scalar cancellation}

\noindent\textbf{Theorem (quadratic scalar Duhamel enclosure).}
Under Eqs.~\eqref{eq:chvi-foundation-full-generator}--
\eqref{eq:chvi-foundation-domain}, define
\begin{equation}
\begin{gathered}
\Gamma_0(s)=\langle\Psi_0,U(s)\Psi_0\rangle,
 \\
 \Gamma_{\rm ref}(s)=e^{-i\bar Es}\cos(\Omega s/2).
\end{gathered}
\end{equation}
Then, for every real $s$,
\begin{equation}
 \boxed{\begin{gathered}
 |\Gamma_0(s)-\Gamma_{\rm ref}(s)|
 \leq\frac{b^2s^2}{2}
 =g^2v_s s^2
 \leq q(s):=g^2v_s^{\rm up}s^2.
 \end{gathered}}
\label{eq:chvi-foundation-quadratic-overlap}
\end{equation}
The actual propagator on the left retains the complete continuum and all
photon sectors. The statement concerns the scalar survival amplitude;
it does not replace the first-order bound on the state-vector difference.

\noindent\emph{Proof.}
Write $\mathcal E(t)=U(t)-R(t)$. Equation~\eqref{eq:chvi-foundation-domain}
permits differentiating $U(t-u)R(u)\Psi_0$ and integrating in the Hilbert
space. With the oriented integral for negative $t$,
\begin{equation}
\begin{gathered}
\mathcal E(t)\Psi_0
   =-i\int_0^t U(t-u)f(u)\,\mathrm du,
 \\
 \|\mathcal E(t)\Psi_0\|\leq b|t|.
\end{gathered}
\label{eq:chvi-foundation-vector-duhamel}
\end{equation}
For $t\geq0$, insert the second identity in
Eq.~\eqref{eq:chvi-foundation-forcing-orthogonality} into the scalar
matrix element to obtain
\begin{equation}
 \langle\Psi_0,\mathcal E(t)\Psi_0\rangle
 =-i\int_0^t
       \langle\Psi_0,\mathcal E(t-u)f(u)\rangle\,\mathrm du.
\end{equation}
Here $\mathcal E(v)^*=\mathcal E(-v)$, so
Cauchy--Schwarz and Eq.~\eqref{eq:chvi-foundation-vector-duhamel} give
\begin{equation}
\begin{split}
 |\langle\Psi_0,\mathcal E(t)\Psi_0\rangle|
 &\leq\int_0^t
       \|\mathcal E(-(t-u))\Psi_0\|\,\|f(u)\|\,\mathrm du\\
 &\leq b^2\int_0^t(t-u)\,\mathrm du
 =\frac{b^2t^2}{2}.
\end{split}
\end{equation}
For negative $t$, self-adjointness gives the complex conjugate scalar
difference at $-t$, with the same modulus. Finally,
Eq.~\eqref{eq:chvi-foundation-preparation} evaluates the reference matrix
element as the displayed cosine. This establishes
Eq.~\eqref{eq:chvi-foundation-quadratic-overlap}. $\square$

The chapter's previous radius was
$\varepsilon_s=|s|\sqrt{2g^2v_s^{\rm up}}$.
The new radius is $q(s)=\varepsilon_s^2/2$. Both bounds remain valid;
one may use $\min\{2,\varepsilon_s,q(s)\}$ at long times. This improvement
does not certify an exact zero, a revival, a phase-clock inverse, or an
extended physical two-arm coupling.

\subsection{Preparation error and timing uncertainty}

\noindent\textbf{Corollary (full-marker preparation error).}
Let $\rho$ be a normalized density operator on the \emph{same complete}
source--field Hilbert space, initially uncorrelated with the route system,
and suppose
\begin{equation}
 T(\rho,|\Psi_0\rangle\langle\Psi_0|)
 :=\tfrac12\|\rho-|\Psi_0\rangle\langle\Psi_0|\|_1
 \leq\delta_{\rm prep}.
\end{equation}
For the unchanged actual generator, set
$\Gamma_\rho(s)=\operatorname{Tr}(\rho U(s))$. Then
\begin{equation}
 |\Gamma_\rho(s)-e^{-i\bar Es}\cos(\Omega s/2)|
 \leq q(s)+2\delta_{\rm prep}.
\label{eq:chvi-foundation-preparation-overlap}
\end{equation}
\emph{Proof.} Trace/operator-norm duality and $\|U(s)\|=1$ bound the
change of amplitude by $2\delta_{\rm prep}$; add
Eq.~\eqref{eq:chvi-foundation-quadratic-overlap}. $\square$

The factor two follows from the stated trace-distance convention.
Closeness of only the source marginal does not supply this full-marker
preparation hypothesis. Trace-distance closeness also does not, by
itself, give a finite energy variance for $\rho$. Neither inference is
needed for Eq.~\eqref{eq:chvi-foundation-preparation-overlap}.

\noindent\textbf{Corollary (centered-energy timing bound).}
Let $s$ be the nominal proper-time difference in the chapter's
dimensionless chart and let the actual difference be $s+\xi$, with
$|\xi|\leq a$. Define the centered amplitude
\begin{equation}
 G_\rho(t)=e^{i\bar Et}\Gamma_\rho(t)
          =\operatorname{Tr}[\rho e^{-i(H_0-\bar E)t}].
\end{equation}
Then
\begin{equation}
 \boxed{\begin{gathered}
 |G_\rho(s+\xi)-\cos(\Omega s/2)|
 \leq q(s)+2\delta_{\rm prep}+\sigma_H a.
 \end{gathered}}
\label{eq:chvi-foundation-centered-jitter}
\end{equation}
A certified upper bound for $\sigma_H$ may replace $\sigma_H$.
If the two route ages have errors $|\xi_j|\leq a_j$, take
$\xi=\xi_1-\xi_2$ and $a=a_1+a_2$; a common equal age error cancels.

\emph{Proof.} For $A=H_0-\bar E$, the spectral theorem and
$|e^{-ix}-1|\leq|x|$ give
\begin{equation}
 \|(e^{-iA\xi}-I)\Psi_0\|
 \leq |\xi|\|A\Psi_0\|=\sigma_H|\xi|.
\end{equation}
Consequently $|G_0(s+\xi)-G_0(s)|\leq\sigma_H|\xi|$.
Compare $G_\rho(s+\xi)$ first with $G_0(s+\xi)$, then with $G_0(s)$,
and finally with the reference cosine. The three errors are bounded by
$2\delta_{\rm prep}$, $\sigma_Ha$ and $q(s)$, respectively. Thus only
the exact preparation's variance is used. $\square$

Centering is an explicit energy-phase assignment, not an instruction to
discard an observable route phase. In raw amplitudes the conclusion is
\begin{equation}
 \left|\Gamma_\rho(s+\xi)
 -e^{-i\bar E(s+\xi)}\cos(\Omega s/2)\right|
 \leq q(s)+2\delta_{\rm prep}+\sigma_Ha.
\label{eq:chvi-foundation-raw-jitter}
\end{equation}
If one instead compares with the \emph{nominal} raw reference
$e^{-i\bar Es}\cos(\Omega s/2)$, add the energy-phase error
\begin{equation}
 2|\cos(\Omega s/2)|\,|\sin(\bar E\xi/2)|
 \leq |\cos(\Omega s/2)|\min\{2,|\bar E|a\}.
\label{eq:chvi-foundation-mean-phase-jitter}
\end{equation}
Thus variance controls the centered timing error and the modulus error;
a numerical absolute-phase bound additionally needs the mean energy or
a specified compensating control phase. No numerical $\bar E$ is inferred
from the gap $\Omega$.

Under a common scalar reassignment $H_0\mapsto H_0+\beta I$ and
$H_{\rm ref}\mapsto H_{\rm ref}+\beta I$, $G_\rho$, $\sigma_H$ and
$q$ are unchanged, whereas $\Gamma_\rho(t)\mapsto
e^{-i\beta t}\Gamma_\rho(t)$. Preserving a complete controlled-route
amplitude requires the corresponding external action assignment
$S_{{\rm ext},j}\mapsto S_{{\rm ext},j}+\hbar\beta\eta_j$ at the
actual route ages. Even energy and counterterms remain counted once.

\subsection{Modulus and phase disks}

\noindent\textbf{Corollary (readable amplitude and phase enclosure).}
Suppose any of the preceding estimates gives a complex disk
$|z-c|\leq r$ for a normalized amplitude $z$. Writing $m=|c|$,
\begin{equation}
 \max\{0,m-r\}\leq |z|\leq\min\{1,m+r\}.
\label{eq:chvi-foundation-modulus-disk}
\end{equation}
If $r<m$, the amplitude is nonzero and its circular phase distance obeys
\begin{equation}
 \operatorname{dist}_{\mathbb S^1}(\arg z,\arg c)
 \leq\arcsin(r/m).
\label{eq:chvi-foundation-phase-disk}
\end{equation}
\emph{Proof.} The modulus bounds are the triangle and reverse-triangle
inequalities, together with $|z|\leq1$. Rotate $c$ to $m>0$.
For $r<m$ the disk lies in the open right half-plane; the distance from
$m$ to the ray through $z$ is $m|\sin\theta|\leq r$ with
$|\theta|<\pi/2$. This gives Eq.~\eqref{eq:chvi-foundation-phase-disk}.
$\square$

For exact timing use $c=e^{-i\bar Es}\cos(\Omega s/2)$ and
$r=q(s)+2\delta_{\rm prep}$. For the centered jitter estimate use
$c=\cos(\Omega s/2)$ and
$r=q(s)+2\delta_{\rm prep}+\sigma_Ha$.
The modulus enclosure is then also an enclosure of the raw amplitude's
modulus, and an ideal recombiner has visibility
$2\sqrt{p_1p_2}|z|$ under its stated assumptions.
For the raw jittered phase, restore $-\bar E(s+\xi)$ explicitly.
If the cosine is negative its reference phase includes $\pi$; at its
zeros, or whenever $r\geq m$, this disk yields no phase enclosure.
A valid local phase branch is therefore confined to an interval on
which the disk excludes zero. Additional beam-splitter, routing and
recombination phases must still be included in an experimental fringe.

